\documentclass[UTF8,fontset=none,11pt,a4paper]{ctexart}
\usepackage{ipara-handbook}
\handbooksetup{DS-07 图的表示与遍历}{408 · 数据结构 DS-07}{Graph · Representation · Coverage}
\begin{document}
\handbooktitle{图的表示与遍历}{关系怎样被表示，并从起点覆盖而不重复失控}{408 · 数据结构 Topic DS-07}

\begin{corebox}{母问题：图的关系如何变成可验证的访问过程？}
图遍历的生成链是
\[
\text{Vertices/Edges}\to\text{Representation}\to\text{Frontier}\to\text{Visited State}\to\text{Expansion}\to\text{Coverage}.
\]
邻接矩阵用空间换常数时间的边查询，邻接表按实际边数保存关系；DFS 与 BFS 都要维护 visited，区别在 frontier 的 LIFO/FIFO 取出规则。遍历只负责覆盖和发现关系，不自动解决最短路、MST 或拓扑目标。
\end{corebox}
\begin{methodbox}{本册定位与 Owner 边界}
\begin{itemize}
\item \kw{Owns：}图、邻接矩阵/邻接表/十字链表/邻接多重表、方向与重边语义、DFS/BFS、visited、连通分量与遍历森林。
\item \kw{Uses：}DS02 的栈/队列端点契约和 DS-B01 frontier 接口；DS08 使用本册表示。
\item \kw{Stops：}最短路、最小生成树、拓扑排序等全局目标归 DS08；网络路由只 Use。
\end{itemize}
\end{methodbox}
\tableofcontents\newpage

\section{表示先于算法}
邻接矩阵 $A[i][j]$ 直接记录边存在性或权值，适合稠密图和频繁边查询，空间 $\Theta(V^2)$。邻接表为每个顶点保存出边列表，空间 $\Theta(V+E)$，遍历一个顶点的邻居只扫描实际出边。无向边需要在两端各登记一次；有向边只登记出发端。平行边和自环是否允许，必须在接口中先固定。

\subsection{十字链表：一条有向弧同时进入出边链与入边链}
普通有向邻接表让枚举顶点 $u$ 的出边很自然，但若要枚举入边，往往还需逆邻接表。十字链表把一条弧记录为
\[
(tail,head,\;tlink,\;hlink,\;info),
\]
其中 $tlink$ 连接同弧尾的下一条出边，$hlink$ 连接同弧头的下一条入边；顶点记录同时保存首条出弧和首条入弧。一个边对象被两条链引用，但只存一份权值和附加信息。

不变量是：每条弧在其尾顶点的出边链中出现一次，也在其头顶点的入边链中出现一次；删除弧必须同时从两条链摘除。它适合同时需要入度、出度和双向邻接扫描的有向图。

\subsection{邻接多重表：无向边是共享对象，而不是两份边副本}
无向邻接表通常把边 $\{u,v\}$ 分别登记在 $u,v$ 的链中。邻接多重表为每条无向边只建一个边节点，保存两个端点 $ivex,jvex$ 以及分别连接两个端点下一条关联边的 $ilink,jlink$。从顶点 $u$ 扫描时，要根据 $u$ 是该边的哪一个端点选择对应 link。

其收益是边级插删与标记只作用于一个共享对象；代价是遍历 link 时必须判断当前端点角色。十字链表面向有向弧的入/出双链，邻接多重表面向无向边的两端关联，二者不能只凭“都有两条链”互换。

\subsection{邻接矩阵还可以承载路径计数}
无权图的邻接矩阵 $A$ 中，$(A^k)_{ij}$ 等于从 $i$ 到 $j$ 的长度恰为 $k$ 的游走条数；这里允许重复顶点或边，因而是 walk count，不等于简单路径数。矩阵乘法的中间求和正好枚举最后一步来自哪个中间顶点。这个结论适合小图或固定步数的计数题，复杂度由矩阵乘法决定，不能拿来替代一般 DFS/BFS 的可达性判定。

邻接矩阵读度时必须先固定方向和对角线语义：无向图一行（或一列）非零项数给度；有向图行和是出度、列和是入度；若有自环，度数是否按 $2$ 计必须与题目约定一致。带权图用 $\infty$ 表示无边时，$0$ 只能表示零权边或自身距离，不能把两者混为同一个哨兵。

\begin{methodbox}{表示选择的题面信号}
只需判断任意两点是否相邻且图较稠密，优先邻接矩阵；按实际边枚举出边，邻接表足够；有向图同时频繁枚举入边与出边，看十字链表；无向图需要对同一边做删除、标记或边级维护，看邻接多重表。选择依据是需要沿哪个关系方向访问，而不是结构名称的新旧。
\end{methodbox}

\section{visited 是覆盖不变量}
遍历的核心不变量是：每个顶点最多被“首次发现”一次，已发现但未展开的顶点位于 frontier，展开后其出边已按表示规则检查。非连通图从一个起点只能覆盖一个可达分量；要得到遍历森林，必须在所有顶点上补一层外循环。
\begin{examplebox}{同一关系，不同 frontier}
从 $0$ 出发，若先把邻居压栈，得到 DFS 的深路径；若把邻居入队，得到 BFS 的层次距离。两者都可以覆盖同一可达集，但只有无权图 BFS 的首次发现层数具有最短边数语义。
\end{examplebox}

\section{DFS、BFS 与遍历森林}
递归 DFS 隐式使用调用栈；显式 DFS 必须约定邻接顺序，否则输出序列不唯一。BFS 使用队列，入队时就标记 visited，不能等出队才标记，否则多个前驱会重复入队。对邻接表，二者时间均为 $O(V+E)$，额外空间为 $O(V)$；矩阵扫描每个顶点所有候选邻居，时间为 $O(V^2)$。

\subsection{四个动词固定遍历状态}
把任何遍历逐行代码都翻译成四个动作：
\[
\boxed{Discover\to Mark\to Enqueue/Push\to Expand}.
\]
Discover 是第一次遇到未访问邻居；Mark 把“已发现”写入 visited；Enqueue/Push 把它加入待处理边界；Expand 是将来取出后扫描邻接表。Mark 必须紧跟 Discover，因为 frontier 中的顶点已经被别的路径“预约”，不能再次加入。

\begin{examplebox}{为什么出队时标记会重复}
设 $0$ 同时连接 $1,2$，且 $1,2$ 都连接 $3$。BFS 展开 $1$ 时发现 $3$ 并入队；若尚未标记，展开 $2$ 时又会把 $3$ 入队。发现时标记则把状态从 `Unseen' 立即改为 `Discovered'，第二条边只被检查，不再制造第二份 frontier 项。
\end{examplebox}

\begin{center}\small
\begin{tabularx}{\linewidth}{L{2.1cm}L{2.4cm}YY}
\toprule
顶点状态 & 所在位置 & 已经保证什么 & 下一动作\\
\midrule
Unseen & 不在 frontier & 尚无已知发现路径 & 被邻边 Discover\\
Discovered & frontier 中 & 已有一条发现路径，visited=true & 等待取出\\
Expanded & frontier 外 & 所有邻边已扫描 & 不再展开\\
\bottomrule
\end{tabularx}
\end{center}

DFS 与 BFS 共享这三态；区别只在 Discovered 顶点按 LIFO 还是 FIFO 取出。遍历森林再增加外层动作：找最小的 Unseen 顶点作为新根，重复同一状态机。

\subsection{图的基本计数：先固定对象，再套公式}
无向简单图的边是无序顶点对，最多有 $\binom{V}{2}$ 条边；有向简单图的弧是有序顶点对，最多有 $V(V-1)$ 条弧。无向图中自环是否允许会改变度数口径：一条普通边对两个端点各贡献 $1$，自环对同一顶点贡献 $2$。因此在无自环无向图中
\[
\sum_{v\in V}\deg(v)=2E,
\]
有向图则分别满足 $\sum \operatorname{indeg}(v)=E$ 与 $\sum \operatorname{outdeg}(v)=E$。

连通性判断必须和边数的前提一起读。无向图若 $E<V-1$ 必不连通；若 $E>V-1$ 必有环；$E=V-1$ 既不能单独推出连通，也不能单独推出无环。连通图的生成树恰有 $V-1$ 条边，非连通图的遍历结果是生成森林。对简单无向图，若
\[
E>\frac{(V-1)(V-2)}2,
\]
则必连通；因为不连通时最多把 $V-1$ 个顶点组成完全图、剩一个孤立点。严格大于不可改成大于等于，等号处存在不连通边界例子。

有向图还要区分弱连通、单向连通与强连通；把方向抹去只能证明基图连通，不能推出强连通。无向图可用 BFS/DFS 染色判断二部性：遇到已染色邻点同色即发现奇环，二部图等价于不含奇数长度回路。以上计数和染色都是遍历的直接扩展，不能把“访问过所有顶点”误写成“已经得到最短路”。

有向强连通简单图至少需要一个有向环覆盖所有顶点，因此 $E\ge V$；上界仍为 $V(V-1)$。若有向边数满足 $E>(V-1)(V-2)$，只能推出把方向抹去后的基图连通，不能推出强连通：所有弧朝同一侧仍可能使任意两点无法互达。题目写“有向图连通”时先确认是弱连通、单向连通还是强连通。

由边数反推无向图分量数时，成功连接两个不同分量的边每条最多让分量数减一，所以 $c\ge\max(V-E,1)$；这个下界可由一棵森林取到。若要让分量尽量多，应把边集中在最少的 $k$ 个顶点中，其中 $k$ 是满足 $\binom{k}{2}\ge E$ 的最小整数，其余顶点孤立，于是 $c\le V-k+1$。这是计数构造，不是遍历顺序结论。

\subsection{BFS 的距离、前驱与路径恢复}
对无权图，从源点 $s$ 入队时置 $dist[s]=0$，首次发现 $v$ 时令
\[
dist[v]=dist[u]+1,\qquad parent[v]=u.
\]
队列的 FIFO 顺序保证顶点按非降距离层被展开；因此首次发现就是最少边数距离。若题目只问可达性，保留 visited 即可；若问具体路径，必须额外保存 parent，并从终点反向回溯后再逆序输出。不可达点保留 $\infty$，不能把“没有前驱”当成距离 $0$。

若每条边权都相同为常数 $c>0$，BFS 仍给出最少边数路径，实际距离为 $c\cdot dist$；若边权不同，层号不再代表代价，应转交 Dijkstra 或其他最短路算法。若题目要求最短路径条数，首次发现时令 $count[v]=count[u]$；再次从同一最短层到达 $v$ 时累加 $count[u]$，但只有满足 $dist[u]+1=dist[v]$ 的边才可计入。这个计数状态和 parent 选择是两个不同接口，不能用“最后一次到达的父节点”代替路径数。

\subsection{遍历实现的等价改写与边界}
递归 DFS 的调用栈记录“当前顶点及其下一条尚未检查的邻边”；显式栈实现若只压入顶点而不记录邻接游标，虽然仍可覆盖，但输出次序可能与递归版本不同。BFS 若用层级循环，可在每层开始记录当前队列长度；若只要距离，则直接传播 $dist$ 更不易漏层。空图、孤立顶点、重边、自环和非连通图必须分别测试：自环只应被 visited 拦截，重边不应生成重复发现，外层循环才负责补齐其他分量。

\subsection{DFS 的颜色状态与有向环}
在有向图中，用白色（未发现）、灰色（已发现但递归尚未返回）、黑色（全部邻边处理完）记录 DFS 状态。沿边遇到灰色顶点即得到后向边，说明当前递归栈形成有向环；遇到黑色顶点不能据此判环，因为它属于已经完成的另一分支。无向图若把父边排除后仍遇到已访问邻点，则得到环证据，但重边和自环必须单独处理。

DFS 生成树的先序取决于邻接访问顺序；因此题目若要求唯一序列，必须给出邻接表顺序或顶点扫描顺序。若只问是否存在环、连通分量数或覆盖集合，序列不唯一不影响结论。把 DFS 的“访问顺序”直接当作拓扑序是错误的，拓扑目标应交给 DS08 的逆后序/Kahn 不变量。

\section{隐式状态图：先定义邻接，再选择 BFS 或 DFS}
网格、密码锁、字符串变换等题目没有显式邻接表，但只要能定义“一个合法操作把状态 $x$ 变成哪些状态”，就已经得到图。建模顺序固定为
\[
\text{State}\to\text{Neighbor Generator}\to\text{Forbidden/Visited}\to\text{Goal}\to\text{Frontier}.
\]
状态编码必须保留未来转移所需信息；visited 应在入队/进入递归时标记，而不是出队后才标记，否则同一状态可能被多个前驱重复加入。要求最少操作次数且每步代价相同时使用 BFS；只问覆盖、分量或自底向上汇总时通常使用 DFS。状态总数有限但无法一次列出时，复杂度写成 $O(|S|+|T|)$，其中 $S$ 是实际发现状态、$T$ 是实际生成转移，不能脱离状态编码只写“$O(n)$”。

\subsection{双向 BFS 的成立条件}
已知单一起点与终点、边可逆或能够生成反向邻居时，可从两端同时维护 frontier，每轮扩展节点数更少的一侧；新状态若已经在另一侧 visited 集合中，路径在此相接。它减少的是常见搜索深度下的状态展开量，不改变最坏状态空间阶。若终点未知、有向边没有可构造的反向操作，或两侧使用了不一致的状态规范化，双向 BFS 的相遇证明不成立。

\subsection{Flood Fill 与岛屿族：变体来自“统计什么”}
二维网格把每个合法单元格视为顶点，四连通或八连通决定邻接。Flood Fill 从种子出发覆盖同色连通分量；若新颜色等于旧颜色，应立即返回，避免反复生成等价转移。visited 可用独立布尔表，也可原地改写网格，但后者会破坏输入，必须写入 API 合同。无论 DFS 还是 BFS，最坏都可能保存 $O(mn)$ 个待处理单元，不能只凭“是岛屿题”宣称队列空间为 $O(\min(m,n))$。

岛屿计数在每次遇到未访问陆地时启动一次遍历；最大面积让一次遍历返回节点数；封闭区域或飞地可先从边界陆地遍历并消去，再统计内部剩余；子岛屿需在覆盖第二张图的分量时累计“所有对应位置在第一张图均为陆地”的合取条件。不同岛屿形状可记录 DFS 方向序列，但必须加入回退记号，使不同分支结构不会编码成同一串；这种相对位移编码通常只消除平移，不会自动消除旋转或镜像，若题目把它们视为同形，还需额外规范化。

\begin{methodbox}{网格题的停止条件}
先固定越界、阻塞、已访问三类拒绝条件，再定义一次访问产生什么不可逆进展。若遍历会修改原网格，调用结束后是否需要恢复必须提前决定；回溯需要撤销“路径选择”，普通连通覆盖通常不撤销 visited。
\end{methodbox}

\section{代码与测试证据}
完整实现位于 \codepath{code/ds07_graph_traversal.hpp}，测试位于 \codepath{code/ds07_graph_traversal_test.cpp}。
\lstinputlisting[language=C++,firstline=12,lastline=92,firstnumber=1,numbers=left,caption={邻接表、BFS、DFS 与可达覆盖}]{30_408/10_数据结构/07_图的表示与遍历/code/ds07_graph_traversal.hpp}
\lstinputlisting[language=C++,firstline=8,lastline=45,firstnumber=1,numbers=left,caption={方向、非连通与空图测试}]{30_408/10_数据结构/07_图的表示与遍历/code/ds07_graph_traversal_test.cpp}
\begin{lstlisting}[language=bash]
clang++ -std=c++17 -Wall -Wextra -Werror -pedantic code/ds07_graph_traversal_test.cpp -o /tmp/ds07_test
/tmp/ds07_test
\end{lstlisting}

\section{边界与概念分离}
\begin{center}\small
\begin{tabularx}{\linewidth}{L{2.5cm}C{1.8cm}C{1.8cm}YY}\toprule
操作 & 邻接表 & 邻接矩阵 & 不变量/前提 & 边界\\\midrule
遍历 & $O(V+E)$ & $O(V^2)$ & 每顶点首次发现一次 & 空图、非连通\\
边查询 & $O(\deg u)$ & $O(1)$ & 方向语义固定 & 重边/自环\\
空间 & $O(V+E)$ & $O(V^2)$ & 只保存约定关系 & 稠密/稀疏\\\bottomrule\end{tabularx}\end{center}
\begin{center}\small
\begin{tabularx}{\linewidth}{L{2.3cm}C{.35cm}L{2.3cm}YY}\toprule
概念 A&$\neq$&概念 B&判据&错因\\\midrule
DFS&$\neq$&BFS&栈 vs 队列&误写距离结论\\
首次发现&$\neq$&出队/出栈&visited 时机&重复入 frontier\\
遍历覆盖&$\neq$&最短路/MST&目标不同&把结构事实当目标证明\\\bottomrule\end{tabularx}\end{center}

\section{做题控制协议}
先标顶点、方向、重边和表示；再确定起点与邻接访问顺序。看到“访问所有可达顶点”选 DFS/BFS，看到“最少边数”才升级到 BFS 距离，看到权值优化转交 DS08。每一步写出 visited 与 frontier 的变化，最后检查是否遗漏非连通分量。
\begin{corebox}{压缩复原}
五问：边存在哪里？当前 frontier 是栈还是队列？何时标记 visited？起点能覆盖哪些分量？题目要覆盖还是要优化？
\end{corebox}
\end{document}
